\chapter{Number fields and their rings of integers}

\section{Integrality}

\begin{definition}{}
    Take \(A \subseteq B \) an inclusion of domains. An element \(x \in B\) is
    said to be \textbf{integral over \(A\)} if it is a root of a monic
    polynomial in \(A[X]\).

    The ring \(B\) is said to be \textbf{integral} over \(A\) if all elements of
    \(B\) are integral.
\end{definition}

\begin{example}{}
    The complex numbers \(\overline{\mathbb{Z}}\) that are integral over
    \(\mathbb{Z}\) are called the \emph{algebraic integers}.
\end{example}

\begin{example}{}
    Let \(d \neq 0, 1\) squarefree. Then
    \[
        \mathbb{Z}[\sqrt{d}] = \{a + b \sqrt{d} : a, b \in \mathbb{Z}\} 
    \]
    is integral over \(\mathbb{Z}\). Indeed if \(x = a + b \sqrt{d}\), then
    \({(x - a)}^2 - b^2 d = 0\) 
\end{example}

\begin{example}{}
    If \(\zeta = e^{\frac{2\pi i}{n}}\) then \(\zeta^{n} = 1\), thus \(\zeta \in
    \overline{\mathbb{Z}}\) 
\end{example}

\begin{proposition}{}
    Let \(A\) be a UFD. Let \(k = \mathcal{Q}{(A)}\) the fraction field of
    \(A\). An element \(x \in K\) is integral over \(A\) if and only if \(x \in
    A\) 
\end{proposition}
\begin{proof}{}
    Write \(x = \frac{p}{q}\), with \(p, q \in A\) and \(\mathrm{mcd}{(p, q)} =
    1\). Then in \(x\) is integral over \(A\), it satisfies an equation
    \[
      x^{d} + a_{d-1} x^{d-1} + \dots + a_{0} = 0 \quad a_{i} \in A \; \forall i
    \]
    Then
    \begin{align*}
       &{\left( \frac{p}{1} \right)} ^{d} + a_{d-1} {\left( \frac{p}{q}
       \right)}^{d-1} + \dots + a_{0} = 0\\
        \implies &p^{d} + a_{d-1} p^{d-1}q + \dots a_{0} q^{d} = 0
    \end{align*}
    hence \(q | p^{d}\) but they are coprime, hence \(q \in  A^{*}\) and thus
    \(x \in A\) 
\end{proof}

Conversely, if \(x \in  A\), it is a root of \(X - x\).

\begin{corollary}{}
    \(\mathbb{Q} \cap \overline{\mathbb{Z}} = \mathbb{Z}\) 
\end{corollary}

\begin{proposition}{}
    Let \(A \subseteq B \) inclusion of domains. Fix \(b_{1}, \dots, b_{n} \in
    B\) not all zero. Set \(P = \sum_{i=1}^{n} A b_{i} \) and take \(b\) such
    that \(bP \subseteq P \).

    We can write
    \[
      b b_{i} = \sum_{i=1}^{n} a_{ij} b_{j} \quad a_{ij} \in A
    \]
    Then introduce \(M = {(a_{ij})} \in \mathcal{M}_n{(A)}\) and let 
    \[
      \rchi_{M} {(T)} = \mathrm{det}{(TI_k - M)}
    \]
    then \(\rchi_M {(b)} = 0\) 
\end{proposition}
\begin{remark}{}
    \(M^{t} \) 
\end{remark}

\begin{proof}{}
    Set \(N = bI_n - M\). We have
    \[
      N \begin{pmatrix}
          b_{1} \\
          \vdots \\
          b_n
      \end{pmatrix} = 0
    \]
    So
    \[
      \mathrm{det} N \begin{pmatrix}
          b_{1} \\
          \vdots \\
          b_n
          \end{pmatrix} = {(\transp{\mathrm{Conj}N})} N \begin{pmatrix}
          b_{1} \\
          \vdots \\
          b_n
      \end{pmatrix} = 0
    \]
    Hence, since there exists \(i_{0}\) with \(b_{i_{0}} \neq 0\), 
    \[
      \mathrm{det}N = 0 = \rchi_M{(b)}
    \]
\end{proof}

\begin{example}{}
    Let \(A = \mathbb{Z}\), \(B = \mathbb{C}\) 
    If we let \(b_{1} = 1, b_{2} = \sqrt[3]{2}, b_{3} = \sqrt[3]{4}\). Then we
    can take \(b = 1 + 3 \sqrt[3]{2} - \sqrt[3]{4}\) 
\end{example}

\begin{corollary}{}
    Let \(A \subseteq B \) inclusion of domains. Take \(x_{1}, \dots, x_{r} \in
    B\). Then \(x_{1}, \dots, x_r\) are integral over \(A\) if and only if
    \(A[x_{1}, \dots, x_r]\) is finitely generated as an \(A\)-module.
\end{corollary}

\begin{proof}\( \)
\begin{itemize}
    \item[\(\implies \)] Assume \(x_{1}, \dots, x_r\) are integral over \(A\).
        By induction, we may and do assume that \(r=1\). Take \(\mu \in A[X]\)
        monic with \(\mu {(x_{1})} = 0\). Take \(f{(x_{1})} = x \in A[x_{1}]\),
        with \(f \in A[X]\). Since \(\mu\) is monic, we can write
        \[
          f = \mu Q + R
        \]
        with \(\mathrm{deg}R < \mathrm{deg} \mu\) and \(Q, R \in A[X]\).
        Finally, \(f{(x_{1})} = R{(x_{1})} \in \sum_{i=1}^{\mathrm{deg}\mu - 1}
        Ax_{1}^{i}\) hence \(A[x_{1}]\) is finitely generated.
    \item[\(\impliedby \)] Assume \(A[x_{1}, \dots, x_r\) is finitely generated
        as an \(A\)-module. Let \(b_{1}, \dots, b_n\) be generators and set \(b
        = x_{i}\). Apply the previous proposition to this setting and get a
        monic \(\rchi \in A[X]\) with \(\rchi {(x_{i})} = 0\). Hence, \(x_{i}\)
        is integral over \(A\).
\end{itemize}
\end{proof}

\begin{corollary}{}
    Let \(A \subseteq B \) inclusion of domains. Let \(\overline{A}\) be the set
    of elements of \(B\) that are integral over \(A\) then \(\overline{A}\) is a
    subring of \(B\) 
\end{corollary}
\begin{proof}{}
    Take \(x, y \in \overline{A}\) and consider \(C = A[x, y]\). Then \(C\) is
    finitely generated as an \(A\)-module. But \(C = A[x, y, x+y, xy]\), hence
    \(x+y, xy \in \overline{A}\) 
\end{proof}
In particular, \(\overline{\mathbb{Z}}\) is a ring.
\begin{corollary}{}
    Take \(A \subseteq B \subseteq C  \) a tower of inclusion of domains. Assume
    \(C\) is integral over \(b\) and \(B\) is integral over \(A\). Then \(C\) is
    integral over \(A\).
\end{corollary}
\begin{proof}{}
    Take \(x \in C\) Let \(f \in B[X]\) monic with \(f{(x)} = 0\) let \(b_{1},
    \dots, b_{n}\) be the coefficients of \(f\) and set \(B' = A[b_{1}, \dots,
    b_{n}] \subseteq B \). Since \(b_{1}, \dots, b_{n}\) are integral over
    \(A\), \(B'\) is a finitely generated \(A\)-module. Since \(x\) is integral
    over \(B'\), \(B'[x]\) is a finitely generated \(B'\)-module. Hence,
    \(B'[x]\) is finitely generated as \(A\)-module, so \(x\) is integral over
    \(A\).
\end{proof}

In particular, if \(\overline{A} = \{x \in B : x \text{ integral over \(A\)}\}
\), then \(\{x \in C : x \text{ integral over \(A\)}\} = \{x \in C : x \text{
integral over \(\overline{A}\)}\} \) 

\begin{definition}{}
    Let \(A \subseteq B \subseteq C  \) domains. 
\begin{enumerate}[label = \arabic*.]
    \item The \textbf{integral closure} of \(A\) in \(B\) is the ring
        \[
          \overline{A} = \{x \in B : x \text{ integral over \(A\)}\} 
        \]
    \item If \(B = \mathcal{Q}{(A)}\), then \(\overline{A}\) is called the
        \textbf{integral closure of \(A\)}.
    \item If \(\overline{A} = A\) we say that \(A\) is \textbf{integrally
        closed} or \textbf{normal}
    \item If \(k\) is a finite extension of \(\mathbb{Q}\) (i.e. a number field)
        then the integral closure of \(\mathbb{Z}\) in \(k\) is called the
        \textbf{ring of integers} of \(k\) and is denoted \(\mathcal{O}_k :=
        i(K) \cap \overline{\mathbb{Z}}\) with \(i : K \hookrightarrow  \mathbb{C}\)
        injective.
\end{enumerate}
\end{definition}

\begin{example}{}
UFDs are integrally closed. In particular, \(\mathcal{O}_{\mathbb{Q}} = \mathbb{Z}\) 
\end{example}

\begin{example}{}
    If \(k = \mathbb{Q}(\sqrt{d)}\) with \(d \neq 0, 1\) squarefree, then
    \(\mathbb{Z}[\sqrt{d}] \subseteq \mathcal{O}_k \). In general (depending on
    the value of \(d\)), \(\mathbb{Z}[\sqrt{d}]\) is not necessarily equal to
    \(\mathcal{O}_k\). If \(d \equiv 1 \mathrm{mod} 4\) then \(\frac{1 +
    \sqrt{d}}{2} \in \mathcal{O}_k \sminus \mathbb{Z}[\sqrt{d}]\) 
\end{example}

\paragraph{The Reference Framework (RF)} is having
\begin{itemize}[label = --]
    \item \(A\) a domain with \(\mathcal{Q}{(A)} = K\) 
    \item \(L / K\) a finite separable extension
    \item \(B \subseteq L \) the integral closure of \(A\)  in \(L\) 
\end{itemize}
In particular, we are interested in the cases with \(A = \mathbb{Z}\) and \(K =
\mathbb{Q}\) 

\section{How to compute rings of integers}
During this section we work with RF.
\begin{proposition}{}
    Take \(x \in L\). Then \(x \in B\) if and only if the minimal polynomial
    \(\mu\) of \(x\) has coefficient in \(A\).
\end{proposition}
\begin{proof}\( \)
\begin{itemize}
    \item[\(\implies \)] Assume \(x \in B\). Take \(f \in A[X]\) monic with
        \(b(x) = 0\). Then \(\mu | f\) in \(K[X]\). So the roots \(z_{1}, \dots,
        z_{n}\) of \(\mu\) are roots of \(f\). Hence, \(z_{1}, \dots, z_{n}\)
        are integral over \(A\). But the coefficients of \(\mu\) are in
        \(A[z_{1}, \dots, z_{n}]\) hence they are integral over \(A\). But they
        are in \(K\) and \(A\) is integrally closed, hence they are in \(A\) and
        \(\mu \in A[X]\) 
    \item[\(\impliedby \)] obvious
\end{itemize}
\end{proof}

\begin{corollary}{}
    Take \(A = \mathbb{Z}\), \(K = \mathbb{Q}\) and \(L = \mathbb{Q}(\sqrt{d})\)
    with \(d \neq 0,1\) squarefree. Then
    \[
      \mathcal{O}_L = \begin{cases}{}
          \mathbb{Z}[sqrt{d}] & d = 2, 3 \text{ mod } 4\\
          \mathbb{Z}{\left(\frac{1 + \sqrt{d}}{2}  \right]} & d \equiv 1 \text{
          mod } 4
      \end{cases}
    \]
\end{corollary}
\begin{proof}{}
    Take \(a + b \sqrt{d} \in L\), with \(a, b \in \mathbb{Q}\) and \(b \neq
    0\). The minimal polynomial is 
    \[
      \mu = {(X - a)}^2 - b^2d = X^2 - 2aX + a^2 - db^2 
    \]
    Hence \(x \in \mathcal{O}_L \iff \mu \in \mathbb{Z}[X] \iff 2a \in
    \mathbb{Z}\) and \(a^2 - db^2 \in \mathbb{Z}\).
    If \(x \in \mathcal{O}_L\), then
    \[
      {(2a)}^2 - 4 {(a^2 - db^2)} = 4db^2 = d{(2b)}^2
    \]
    Since \(d\) is squarefree, \(2b \in \mathbb{Z}\). Write \(a = \frac{m}{z}\),
    \(b = \frac{n}{z}\) with \(m, n \in \mathbb{Z}\). We conclude in the two
    cases with the thesis.
\end{proof}

We introduce some algebraic tools. \(L / K\) is finite separable. By the
primitive element theorem we know \(L = K(\alpha)\), with \(\alpha \in L\) and
\(L = K{(\alpha)} \cong K[X] / {(\mu)}\) where \(\mu\) is the minimal polynomial
of \(\alpha\).

Take \(L^s\) a separable closure of \(L\). \(\mathrm{Hom}_K {(L, L^s)} =
\{\sigma : L \hookrightarrow L^s : \sigma|_K = \mathrm{id}_K\} \cong
\{\text{roots of \(\mu\) in \(L^s\)}\}\). In particular, \(|\mathrm{Hom}_K{(L,
L^s)}| = [L : K]\). Coming back to the question of finding out if some
elements of \(L\) are integral or not over \(A\), we know how to compute
polynomials that vanish at some \(x \in L\):
\[
  \begin{align*}
      m_x: L &\longrightarrow L \\
      y &\longmapsto m_x(y) = xy
  \end{align*}
\]
and \(\rchi_{m_x} {(x)} = 0\) 

\begin{definition}{}
    \(\rchi_{m_x} =: \rchi_x\) is called the characteristic polynomial of \(x\).
    \[
      \rchi_x = X^{n} - a_{n-1} X^{n-1} + \dots + {(-1)}^{n} a_{0} 
    \]
    and \(a_{n-1} = \mathrm{Tr}_{L / K} {(x)} \) the trace of \(x\) , \(a_{0} =
    N_{L / K}{(x)}\) the norm of \(x\).
\end{definition}

\begin{proposition}{}
    We have: \(\mu_x\) minimal polynomial of \(x\). Then
    \[
        \rchi_x = {(\mu_x)}^{[L  K(x)]} = \prod_{\sigma \in  \mathrm{Hom}_K{(L,
        L^s)}} {(X - \sigma {(x)})}
    \]
    and 
    \[
      \mathrm{Tr}_{L / K} {(x)} = \sum_{\sigma \in \mathrm{Hom}_{K} {(L,
      L^s)}} \sigma {(x)} \quad; \quad N_{L / K} {(x)} = \prod_{\sigma \in
  \mathrm{Hom}_K{(L, L^s)}} \sigma {(x)}
    \]
\end{proposition}

\begin{proof}{}
    %TODO integrare con foto
    \[
      \begin{pmatrix}
          0 & 0 & \dots & 0 & -a_{n-1} \\
          1 & 0 & \dots & 0 & -a_{n-2} \\
          0 & 1 & \dots & \vdots & \vdots \\
          \vdots & \vdots & \ddots & 0 & \vdots \\
          0 & 0 & \dots & 1 & -a_{0}
      \end{pmatrix}
    \]
    where \(\mu = X^{n} + a_{n-1} X^{n-1} + \dots + a_{0}\). Since the
    characteristic poly of tthis matrix is \(\mu\), we get
    \[
      \rchi_x = \mu_x^{v} = \mu_x^{[L : K{(x)}}
    \]
    We know:
    \begin{align*}
        \mu_x &= \prod_{\tau \in \mathrm{Hom}_K{(K{(x)}, L^{s})}} {(X -
      \tau{(x)})}\\
            \rchi_x &= \prod_{\tau \in \mathrm{Hom}_K{(K{(x)}, L^{s})}} {(X -
            \tau {(x)})}^{[L : K{(x)}]}
    \end{align*}
    We have to check that, for each \(\tau \in \mathrm{Hom}_K {(K{(x)},
    L^{s})}\), there exists exactly \([L : K{(x)}]\) elements \(\sigma \in
    \mathrm{Hom}_K{(L, L^{s})}\) with \(\sigma|_{K{(x)}} = \tau\).

    Then, by PET, \(L = K{(x, y)}\) with some \(y \in L\).
    Let \(\mu_y\) be the minimal polynomial of \(y\) over \(K{(x)}\).
    Fixing \(\tau \in  \mathrm{Hom}_K{(K{(x)}, L^{s})}\), 
    \[
    \mathrm{Hom}_{K{(x)}} {(L, L^{s}})} = \{\sigma : L \hookrightarrow L^{s} :
\sigma|_{K{(x)}} = \tau\} \text{ has } [L : K{(x)}] = \mathrm{deg}\mu_y \text{
elements }
    \]

\end{proof}

\begin{corollary}{}
    Take \(x \in L\). Then \(x \in B\) if and only if \(\rchi_x \in A[X]\). In
    particular, if \(x \in B\), then \(\mathrm{Tr}_{L / K} {(x)} \in A\) and
    \(N_{L / K} {(x)} \in  A\) 
\end{corollary}
\begin{proof}{}
    
\end{proof}
Fix \(\mathbf{x} = {(x_{1}, \dots, x_{n})}\) a \(K\)-basis of \(L\). Given \(x =
\sum_{i=1}^{n} d_{i} x_{i} \), \(d_{i} \in K\), I want to express in terms of
the \(d_{i}\)'s whether \(x \in B\) or \(x \not\in B\).

\begin{lemma}{}
    \(\mathcal{Q}{(B)} = L\). Better, for each \(x \in L\), there exists \(y \in
    A \sminus \{0\} \) such that \(yx \in B\).
\end{lemma}
\begin{proof}{}
    Take \(\mu_x \in K[X]\) the minimal polynomial of \(x\). Write \(a_{i} =
    b_{i} / c_{i}\), with \(b_{i}, c_{i} \in A\). Set \(c := {(c_{0}, \dots,
    c_{m-1} )}\) and \(z := cx\). Then
    \[
      {\left( \frac{z}{c} \right)}^{m} + a_{m-1} {\left( \frac{z}{c} \right)}
      ^{m-1} + \dots + a_{0} = 0
    \]
    by multiplying everything times \(c^{m}\) we find \(z \in B\), hence \(x =
    \frac{z}{c}\) with \(z \in B\) and \(c \in A\).
\end{proof}

Up to multiplying \(\mathbf{x} \) by some well-chosen \(a \in A\), we may assume
\(x_{1}, \dots, x_{n} \in B\). We come back to 
\[
  x = \sum_{i=1}^{n} \lambda_{i} x_{i}, \quad \lambda_{i} \in K 
\]
If \(x \in B\), then \(\forall i\), \(x x_{i} \in B_n\) and \(\mathrm{Tr}_{L /
K} {(x x_{i})} = \sum_{j=1}^{n} \lambda_{i} \mathrm{Tr}_{L / K} {(x_{i} x_{j})}
\in A\).

If we set \(M_{\mathbf{x} } = {\left( \mathrm{Tr}_{L / K} {(x_{i}x_{j})}
\right)} \in \mathcal{M}_n{(A)} \), then \(M_{\mathbf{x} } \begin{pmatrix}
    \lambda_{1} \\
    \lambda_{2} \\
    \vdots \\
    \lambda_n
\end{pmatrix} \in A^{n}\) and hence its determinant is in \(A^{n}\).

\begin{proposition}{}
    If \(d{(\mathbf{x} )} = \mathrm{det}{(M_{\mathbf{x} } )} =
    \mathrm{det}{\left( \mathrm{Tr}_{L / K} {(x_{i}x_{j})} \right)} \) and
    \(d{(\mathbf{x} )} B \subseteq \sum_{i=1}^{n} Ax_{i}  \) 
\end{proposition}

\begin{definition}{discriminant}
    d\(\mathbf{x} \) is called the \textbf{discriminant} of \(\mathbf{x} \) 
\end{definition}

\paragraph{Computation of the discriminant}
\begin{align*}{}
    d{(\mathbf{x} )} &= \mathrm{det}{\left( \mathrm{Tr}_{L / K} {(x_{i}x_{j})}
    \right)} = \\
                     &= \mathrm{det} {\left( \sum_{\sigma \in
                     \mathrm{Hom}_K{(L, L^{s})}} \sigma{(x_{i})}\sigma{(x_{j})}
             \right)} =\\
                     &= {\left( \mathrm{det}{(\sigma {(x_{j})})}_{\sigma, j}
                     \right)}^2
\end{align*}
% If \(\mathbf{x}  = {(1, x, x^2, \dots, x^{n-1})}\) then

\subsection{Theoretical algorithm to compute rings of integer}

Take \(A = \mathbb{Z}\), \(K = \mathbb{Q}\), \(L\) a number field and \(B =
\mathcal{O}_L\), then
\begin{enumerate}[label = \arabic*)]
    \item Find a \(\mathbb{Q}\)-basis \(\mathbf{x} = {(x_{1}, \dots, x_{n})}\)
        of \(L\) contained in \(B\) e.g. \(\mathbf{x} = {(1, x, x^2, \dots,
        x^{n-1})}\) with \(x \in B\) a primitive element for \(L / \mathbb{Q}\) 
    \item Compute \(d{(\mathbf{x} )}\) 
    \item Say that
        \[
          \bigoplus_{i} \mathbb{Z} x_{i} \subseteq B \subseteq
          \frac{1}{d{(\mathbf{x} )}} \bigoplus \mathbb{Z}x_{i}
        \]
    \item Compute representation \(y_{1}, \dots, y_{n}\) of the quotient
        \[
          {\left( \frac{1}{d{(\mathbf{x} )}} \bigoplus \mathbb{Z} x_{i}\right)}
          / {\left( \bigoplus \mathbb{Z} x_{i}\right)} 
        \]
\end{enumerate}

\begin{proposition}{}
    In the RF. Assume \(A\) is a PID. Then
\begin{enumerate}[label = \arabic*.]
    \item There exists a \(K\)-basis \(\mathbf{x} = {(x_{1}, \dots, x_{n})}\) of
        \(L\) contained in \(B\) such that
        \[
          B = \bigoplus_{i} A x_{i}
        \]
        such a basis \(\mathbf{x} \) is called an \textbf{integral basis} of \(B\) 
    \item Assume further that \(A = \mathbb{Z}\). Then \(d{(\mathbf{x} )}\) is
        independant of the choice of an integral basis \(\mathbf{x} \) and if
        \(\mathbf{y} \) is any \(K\)-basis of \(L\) contained in \(B\), then 
        \[
            d{(\mathbf{y} )} = [ B : \Gamma_\mathbf{y} ]^{2} \cdot d{(\mathbf{x} )}
        \]
\end{enumerate}
\end{proposition}
\begin{remark}{}
    If \(A = \mathbb{Z}\) and \(K = \mathbb{Q}\). Say you have a
    \(\mathbb{Q}\)-basis \(\mathbf{y} \) of \(L\) contained in \(\mathcal{O}_L\)
    such that \(d{(\mathbf{y} )}\) is squarefree. Then \(\mathcal{O}_L =
    \bigoplus \mathbb{Z} y_{i}\) 
\end{remark}

\begin{proof}{}
\begin{enumerate}[label = \arabic*)]
    \item Take \(\mathbf{y} \) a \(K\)-basis of \(L\) contained in \(B\). We
        know
        \[
          \bigoplus_i A y_{i} \subseteq B \subseteq \frac{1}{d{(\mathbf{y} )}}
          \bigoplus_i A y_{i}
        \]
        By classication of finitely generated modules over PIDs, \(B\) is a free
        \(A\)-module of rank \([L : K]\). Hence, there exists \(\mathbf{x}  =
        {(x_{1}, \dots, x_{n})}\) such that \(B = \bigoplus_i A x_{i}\).
        Necessarily, \(\mathbf{x} \) is also a \(K\)-basis of \(L\) 
    \item Let \(\mathbf{x} \) be an integral basis of \(B = \mathcal{O}_L\). Let
        \(\mathbf{y} \) be any \(\mathbb{Q}\)-basis of \(L\) contained in
        \(\mathcal{O}_L\). We know that \(d{(\mathbf{x})}\) (resp. d\(\mathbf{y}
        \)) is the determinant of the trace form \(\mathrm{Tr}_{L / K} \) on the
        basis \(\mathbf{x} \) (resp. \(\mathbf{y} \)). If \(P\) is the
        transition matrix from \(\mathbf{x} \) to \(\mathbf{y} \), then
        \(d{(\mathbf{y} )} = d{(\mathbf{x} )} {\left( \mathrm{det} P \right)}^2
        \). By Smith's nor form theorem (théorème de la base adaptée),
        \(|\mathrm{det}P| = [ \bigoplus_{i} \mathbb{Z} x_{i} : \bigoplus_{i}
        \mathbb{Z}y_{i} ] = [\mathcal{O}_L : \Gamma_\mathbf{y} ]\). We conclude
        with the equivalence in the thesis. Moreover, if \(\mathbf{y} \) was an
        integral basis, then \(\Gamma_\mathbf{y} = \mathcal{O}_L\) and thus
        \(d{(\mathbf{y} )} = d{(\mathbf{x} )}\) 
\end{enumerate}

\end{proof}

\begin{definition}{}
    Let \(L\) be a number field. The discriminant of an integral basis of \(L\)
    is called \textbf{discriminant} of \(L\) and it's noted \(\Delta_L\) 
\end{definition}

\section{Algebraic properties}

Rings of integers do not satisfy many properties. They are not PID generally for
instance.
\begin{definition}{Dedekind domains}
    A domain \(A\) is said to be a \textbf{Dedekind domain} if it is Noeherian,
    integrally closed and every non-zero prime ideal in \(A\) is maximal.
    Equivalently, it's a Noetherian domain with Krull dimension at most 1.
\end{definition}

\begin{example}{}
    PIDs are Dedekind domains. UFDs are not all Dedekind domains. For instance,
    in \(A = \mathbb{C}[x, y] \) \((x)\) is not maximal.
\end{example}

\begin{theorem}{}
    Working in the Framework RF (even though separability of \(L / K\) is not
    needed). Assume \(A\) is a Dedekind domain. Then \(B\) is also a Dedekind
    domain.
\end{theorem}
\begin{remark}{}
    In particular, rings of integers in number fields are Dedekind domains.
\end{remark}
\begin{proof}{}
    (We use the assumption of separability of \(L / K\)). We already know that
    \(B\) is integrally closed. We know that, if \(\mathbf{x} = {(x_{1}, \dots,
    x_{n})}\) is a \(K\)-basis of \(L\) contained in \(B\) and \(B \subseteq
    \frac{1}{d{(\mathbf{x} )}} \bigoplus A x_{i} \) which is a finitely
    generated \(A\)-module, hence noetherian. It follows that \(B\) is a
    noetherian ring.

    Take a non-zero prime ideal \(\mathfrak{p}\) in \(B\). Set \(\mathfrak{q} =
    \mathfrak{p}^{c} = \mathfrak{p} \cap A\). Then \(\mathfrak{q}\) is a prime
    ideal in \(A\). Moreover, if one takes \(x \in \mathfrak{p} \sminus \{0\}
    \), then \(x\) satisfies some equation
    \[
      x^{n} + a_{n-1} x^{n-1} + \dots + a_{0} = 0 \quad a_{i} \in  A, \; a_{0}
      \neq 0
    \]
    Hence 
    \[
      a_{0} = -x^{n} - a_{n-1} x^{n-1} - \dots - a_{1} x \in  \mathfrak{p} \cap
      A = \mathfrak{q}
    \]
    thus \(\mathfrak{q}\) is non-zero, hence maximal. 
    \(B / \mathfrak{p}\) is a domain which is also a finite-dimensional \({(A /
    \mathfrak{q})}\)-vector space (\(B\) is finitely generated \(A\)-module).
    Hence, \(B / \mathfrak{p}\) is a field and \(\mathfrak{p}\) is maximal.
\end{proof}

\subsection{Fractional ideals}
\begin{definition}{}
    Let \(A\) be a domain with \(\mathcal{Q}{(A)} = K\). \(A\) fractional ideal
    \(I\) in \(K\) is an \(A\)-submodule of \(K\) such that there exists \(a \in
    A \sminus \{0\} \) satisfying \(aI \subseteq A \) 
\end{definition}

\begin{example}{}
    If \(A\) is a PID, fractional ideals are the \(xA\) with \(x \in K \sminus
    \{0\} \). In particular if \(A\) is not a field, \(K\) is not a fractional
    ideal.
\end{example}

\begin{theorem}{}
    Let \(A\) be a Dedekind domain. All non-zero prime ideals in \(A\) are
    invertible.
\end{theorem}
Take \(\mathfrak{p}\) a non-zero prime ideal in \(A\). If \(J\) is a frctional
ideal such that \(\mathfrak{p}J \subseteq A \), then we define
\[
  J \subseteq \{ x \in  K : x \mathfrak{p} \subseteq A \} =: \mathfrak{p}^{-1}  
\]
Note that if \(y \in  \mathfrak{p} \sminus \{0\} \), then \(y\mathfrak{p}^{-1}
\subseteq A \), hence \(\mathfrak{p}^{-1}\) is a fractional ideal.

The theorem will then follow from the following 
\begin{proposition}{}
    Let \(A\) be a Dedekind domain, \(\mathfrak{p}\) a non-zero prime ideal in
    \(A\) and \(I\) any non-zero ideal in \(A\) and \(I\) any non-zero ideal in
    \(A\). Then \(I \subsetneq \mathfrak{p}^{-1} I \). In particular,
    \(\mathfrak{p}^{-1} \mathfrak{p} = A\) 
\end{proposition}
\begin{remark}{}
    In particular, \(\mathfrak{p} \subsetneq \mathfrak{p}^{-1} \mathfrak{p}
    \subseteq A  \) since \(\mathfrak{p}\) is maximal, \(\mathfrak{p}^{-1}
    \mathfrak{p} = A\) 
\end{remark}
\begin{proof}{}
    \emph{Step 1:}
    We prove that every nonzero ideal \(I \subseteq A \) contains a product of
    nonzero prime ideals. Assume by contradiction this is wrong. Consider
    \(\mathcal{E}\) a set of nonzero ideals \(I \subseteq A \) that do not
    contain a product of nonzero prime ideals.
    By assumption \(\mathcal{E} \neq \varnothing\) and using Notherianity, it
    has a maximal element \(I_{0}\). \(I_{0}\) is not prime, so we can find \(a,
    b \in A \sminus I_{0}\) such that \(ab \in I_{0}\). Then \(I_{0} + {(a)}\)
    and \(I_{0} + {(b)}\) both contain product of nonzero prime ideals.
    Hence, \({(I_{0} + {(a)})}{(I_{0} + {(b)})} \subseteq I_{0} \) contains a
    product of nonzero prime ideals, which leads to a contradiction.

    \emph{Step 2:} We prove that \(A \subsetneq \mathfrak{p}^{-1} \).
    \(\supseteq  \) is obvious. Take now \(a inn \mathfrak{p} \sminus \{0\} \);
    by \emph{Step 1} we can write \(\mathfrak{p}_1 \dots \mathfrak{p}_r
    \subseteq {(a)} \subseteq  \mathfrak{p} \) and we suppose \(r\) as small as
    possible, with all \(\mathfrak{p}_i\) prime ideals. There exists \(i_{0}\)
    such that \(\mathfrak{p}_{i_{0}} \subseteq \mathfrak{p} \) and we may assume
    \(i_{0} = r\). Since \(\mathfrak{p}_r\) is maximal, \(\mathfrak{p}_r =
    \mathfrak{p}\). Moreover, since \(r\) is as small as possible, we can find
    \(b \in  \mathfrak{p}_1 \dots \mathfrak{p}_{r-1} \sminus {(a)} \). Then
    \(a^{-1}b \not\in A\) and \(\forall x \in \mathfrak{p}\), \(a^{-1}b x \in
    a^{-1} \mathfrak{p}_1 \dots \mathfrak{p}_{r-1} \mathfrak{p} =
    a^{-1}\mathfrak{p}_1 \dots \mathfrak{p}_r \subseteq a^{-1} {(a)} =  A\).
    Hence, \(a^{-1}b \in  \mathfrak{p}^{-1}\).

    \emph{Step 3:} We prove the proposition \(\mathfrak{p}^{-1} I \supsetneq I
    \) by contradiction. Assume \(\mathfrak{p}^{-1} I = I\) and take \(x \in
    \mathfrak{p}^{-1}\). Then \(x \in \mathfrak{p}^{-1}\). Then \(xI \subseteq I
    \) and by Notherianity of \(A\), \(I\) is finitely generated. By a previous
    proposition, this allows to find a monic \(\rchi \in  A[X]\) such that
    \(\rchi {(x)} = 0\). But \(A\) is integrally closed, hence \(x \in A\).
    Hence \(\mathfrak{p}^{-1} = A\) which contradicts \emph{Step 2}.
\end{proof}

\subsection{Factorization of ideals}
\begin{theorem}{}
    Let \(A\) be a Dedekind domain. Then every nonzero ideal of \(A\) can be
    written as a product of nonzero prime ideals, in a unique way up to the
    order.
\end{theorem}
\begin{proof}{}
    \emph{Existence:} Let \(\mathcal{E}\) be the family of ideals in \(A\) other
    than \({(0)}\) and \({(1)}\) that are not product of prime ideals. By
    contradiction, assume \(\mathcal{E} \neq 0\). By Notherianity, we can find a
    maximal element \(I_{0} \in \mathcal{E}\). Since \(I_{0} \neq A\) we can
    find a maximal ideal \(\mathfrak{p}\) of \(A\) containing \(I_{0}\). Set \(J
    = \mathfrak{p}^{-1} I_{0}\).

    Then
    \[
      i_{0} \subsetneq \mathfrak{p}^{-1}I_{0} \subseteq \mathfrak{p}^{-1}
      \mathfrak{p} = A 
    \]
    So \(\mathfrak{p}^{-1}I_{0}\) is an ideal of \(A\) not in \(\mathcal{E}\).
    So \(\mathfrak{p}^{-1} I_{0} = \mathfrak{p}_1 \dots \mathfrak{p}_r\) for
    some prime ideals \(\mathfrak{p}_1, \dots, \mathfrak{p}_r\). We conclude
    \(I_{0} = \mathfrak{p} \mathfrak{p}^{-1} I_{0} = \mathfrak{p} \mathfrak{p}_1
    \dots \mathfrak{p}_r\) 

\end{proof}



\section{Ramification}

\begin{definition}{}
    In the RF with \(A\) Dedekind, take \(\mathfrak{p}\) a prime ideal in \(A\).
    Write the factorization of \(\mathfrak{p}B = p_{1}^{e_{1}} \dots p_g
    ^{e_g}\).
\begin{enumerate}[label = \arabic*.]
    \item 
\end{enumerate}

\end{definition}


\begin{proposition}{}
    Almost all prime ideals of \(A\) are unramified in \(L\).  
\end{proposition}
\begin{proof}{}
    Let \(\alpha \in
    B\) a primitive element for \(L / K\) and take \(\mu\) its minimal
    polynomial. Set
\end{proof}

\begin{theorem}{}
    Consider the \emph{discriminant ideal} \(\mathfrak{d}_{B / A} \) of \(A\)
    generated by the discriminants of all \(K\)-bases of \(L\) contained in
    \(B\).

    A prime ideal \(\mathfrak{p}\) of \(A\) is ramified in \(L\) \(\iff\)
    \(\mathfrak{p} | \mathfrak{d}_{B / A} \) 
\end{theorem}



